\section{统一实验矩阵、优先级与执行顺序}

\begin{table}[htbp]
\centering
\caption{补强实验统一矩阵。列为四类指标（证据召回、端到端质量、时延/成本、相关结构）；行为实验单元（选择器 $\times$ 配置 $\times$ 数据集 $\times$ 阅读器）。P0 必做、P1 强烈建议、P2 视资源；``本地''指纯 CPU 容器计算，``API''指阅读器调用。}
\label{tab:matrix}
\small
\begin{tabularx}{\textwidth}{@{}lXcccl@{}}
\toprule
实验 & 范围 & 指标 & 依赖 & 优先级 & 资源 \\
\midrule
C1 消融跨数据集 & 5 配置 $\times$ 3 池，$r{=}4$ & 召回+区间 & 无 & P1 & 本地 \\
C2 配置$\times$预算矩阵 & 11 配置 $\times$ 4 预算 $\times$ 3 池 & 召回+区间 & 无 & P1 & 本地 \\
C3 frontier 重绘 & C2 + 既有门测量 & 三轴 & C2 & P1 & 本地 \\
A0 恢复 $n{=}100$ & 7 方法，QASPER，$r{=}4$ & EM/F1/保留 & API 配额 & P0 & API \\
A1 $n{=}300$ & 7 方法，QASPER，$r{=}4$ & EM/F1/保留 & A0 & P1 & API \\
A2 跨数据集 e2e & 4 方法 $\times$ 2 池，$n{=}100$ & EM/F1/保留 & A0 & P2 & API \\
B1 多阅读器 & 7 方法 $\times$ 3 阅读器，$n{=}100$ & EM/F1 & A0 & P2 & API+GPU \\
B2 位置剖面重测 & 5 位置 $\times$ 3 阅读器，$n{=}50$ & F1 曲线 & 无 & P2 & API+GPU \\
B3 相关研究扩展 & 3 阅读器 $\times$ 方法对 & Pearson & B1 & P2 & 复用 \\
C4 预算自适应（可选） & IDF 分桶调 $B$，3 池 & 召回 & C1 & 可选 & 本地 \\
\bottomrule
\end{tabularx}
\end{table}

执行顺序按依赖与成本排成三个波次。第一波（本周内可完成，除 A0 外全部零 API）：C1 与 C2 并行——两者改同一脚本的不同循环，合并为一次提交；C3 在其落盘后重绘。第一波完成后修订稿即可获得跨数据集消融表与统一矩阵表两张新表与一张升级图，``premature''的第二构成（消融未连成矩阵）即被消解。第二波：A0（API 配额恢复即挂机）与 A1 串行，A1 在 A0 的断点机制上参数化扩展；完成后端到端主表换血，``premature''的第一构成消解过半。第三波（视审稿窗口与预算）：A2 与 B1/B2/B3，其中 B 系列共享 GPU 时租与阅读器端点配置，一次搭好后三项连续执行；完成后第三构成（单点依赖）消解，且阅读器侧位置剖面从``借用''升格为``自测''。

总预算量级汇总：API 调用约 $700+1{,}400+800+2{,}100+750\approx 5{,}750$ 次阅读器问答（A0 至 B2 全做的情形；每问输入中位约 1{,}500 token、输出上限 96 token），本地计算合计不超过单容器一天，GPU 时租为两个本地阅读器服务时长（B 系列内复用）。若预算受限，裁剪顺序为：先砍 C4，再砍 A2（保留 B1 的跨阅读器主张，放弃跨数据集端到端），最后把 B1 的 API 档阅读器从二减一（保留本地开源档与既有 GLM 档，仍有三阅读器中的两个新证据点）。
